\section{A. 分析学（Analysis）}

\subsection{A1. 从黎曼积分到勒贝格积分}

\subsubsection{黎曼积分的本质与局限}

黎曼积分的核心思想：将定义域 $[a,b]$ 分割为 $n$ 个小区间，在每个小区间上取函数值乘以区间长度求和，取极限：
\[
\int_a^b f(x)\,dx = \lim_{\max \Delta x_i \to 0} \sum_{i=1}^n f(\xi_i)\,\Delta x_i
\]

\textbf{局限性：}
\begin{itemize}
  \item 黎曼可积的充要条件：$f$ 有界且不连续点集合的勒贝格测度为零（即“几乎处处连续”）
  \item 狄利克雷函数（有理数处为 1，无理数处为 0）黎曼不可积
  \item 黎曼积分对极限运算不友好：$f_n \to f$ 逐点收敛时，$\int f_n$ 不一定收敛到 $\int f$
\end{itemize}

\subsubsection{勒贝格积分的构造}

核心思想转变：\textbf{不再分割定义域，而是分割值域。}

对于函数 $f: X \to \mathbb{R}$，考虑值域中的水平集：
\[
E_y = \{x : f(x) > y\}
\]

勒贝格积分定义为：
\[
\int f\,d\mu = \int_0^\infty \mu(\{x: f(x) > t\})\,dt
\]

构造步骤：
\begin{enumerate}
  \item \textbf{简单函数：} $s(x) = \sum_{i=1}^n a_i \mathbf{1}_{A_i}(x)$，其中 $A_i$ 是可测集
  \item 定义 $\int s\,d\mu = \sum a_i\,\mu(A_i)$
  \item 对非负可测函数 $f$：$\int f\,d\mu = \sup\left\{\int s\,d\mu : 0 \leq s \leq f,\; s \text{ 简单}\right\}$
  \item 一般函数：$f = f^+ - f^-$，分别积分
\end{enumerate}

\subsubsection{三大收敛定理}

\begin{center}
\begin{tabular}{L{3.5cm}L{5cm}L{4cm}}
\toprule
\textbf{定理} & \textbf{条件} & \textbf{结论} \\
\midrule
单调收敛定理 (MCT) & $0 \leq f_1 \leq f_2 \leq \cdots$，$f_n \to f$ & $\int f_n \to \int f$ \\
控制收敛定理 (DCT) & $f_n \to f$ a.e.，$\exists$ 可积 $g$ 使 $|f_n| \leq g$ & $\int f_n \to \int f$ \\
Fatou 引理 & $f_n \geq 0$ & $\int \liminf f_n \leq \liminf \int f_n$ \\
\bottomrule
\end{tabular}
\end{center}

在 CG 中：渲染方程的蒙特卡洛估计收敛性证明依赖 DCT。

\subsubsection{测度论基础}

\textbf{$\sigma$-代数：} 集合 $X$ 上的 $\sigma$-代数 $\mathcal{F}$ 是 $X$ 的子集族，满足：
\begin{itemize}
  \item $X \in \mathcal{F}$
  \item 对补集封闭：$A \in \mathcal{F} \Rightarrow A^c \in \mathcal{F}$
  \item 对可数并封闭：$A_1, A_2, \ldots \in \mathcal{F} \Rightarrow \bigcup A_i \in \mathcal{F}$
\end{itemize}

\textbf{为什么是“可数”并，而不是“任意”并？} 这是为了与下面测度的\textbf{可数可加}相匹配。
$\mathbb{R}$ 是可分的（存在可数拓扑基），每个开集都是可数个基开集之并，因此 Borel $\sigma$-代数
完全由\textbf{可数}层级的操作生成（超限层级 $\Sigma^0_\alpha$）。若改成对\emph{任意}并封闭，
$\mathcal{F}$ 立刻退化为幂集 $2^X$——而 Vitali 集证明：在幂集上无法同时拥有可数可加与平移不变的测度。
反过来说，$\sigma$-代数的“可数”二字，正是为可数可加腾出的逻辑空间。

\textbf{测度：} $\mu: \mathcal{F} \to [0, +\infty]$，满足：
\begin{itemize}
  \item $\mu(\emptyset) = 0$
  \item \textbf{可数可加性}：若 $\{A_n\}_{n=1}^\infty \subset \mathcal{F}$ 两两不交，则
  \[
  \mu\Big(\bigcup_{n=1}^\infty A_n\Big) = \sum_{n=1}^\infty \mu(A_n)
  \]
\end{itemize}

\textbf{可数可加的必要性——为什么恰好是“可数”？}\\
可数可加是\textbf{唯一同时满足两个要求的强度}：强到足以支撑分析，弱到不排除连续测度。

\textit{(1) 有限可加太弱：丢失全部极限工具。} 有限可加\emph{不}蕴含可数可加，有三个经典反例：

\begin{center}
\begin{tabular}{L{2.8cm}L{5.6cm}L{4.4cm}}
\toprule
\textbf{集函数} & \textbf{定义} & \textbf{失效方式} \\
\midrule
自然密度 & $\mu(A)=\lim\limits_{n}\frac{|A\cap\{1,\dots,n\}|}{n}$（极限存在时） &
$\mu(\{k\})=0$ 但 $\mu(\mathbb{N})=1$ \\
\addlinespace
退化测度 & $\mu(A)=0$（$A$ 有限），$\mu(A)=\infty$（$A$ 无限） &
$\mathbb{N}=\bigsqcup_k\{k\}$ 却 $\mu(\mathbb{N})=\infty \neq \sum_k 0$ \\
\addlinespace
Banach 测度 & $\mathbb{R}$ 上 Lebesgue 测度的有限可加扩张 &
三个条件全满足，但无任何收敛定理 \\
\bottomrule
\end{tabular}
\end{center}

Banach 测度说明：\emph{只}要求有限可加时，可以做得“漂亮”——但代价是\textbf{一整套定理同时失效}：

\begin{center}
\begin{tabular}{L{5.2cm}L{7.8cm}}
\toprule
\textbf{失去的东西} & \textbf{后果} \\
\midrule
MCT / Fatou / DCT & 积分与极限不能交换 $\to$ 无法做分析 \\
$L^p$ 完备性（Riesz--Fischer） & 泛函分析地基消失 \\
Fubini--Tonelli & 累次积分不成立 \\
Radon--Nikodym、Lebesgue 分解 & 无“密度函数”，条件期望无法定义（$\to$ C4） \\
Carathéodory 扩张的\textbf{唯一性} & 由区间长度构造 Lebesgue 测度时失去确定性 \\
Kolmogorov 相容性定理 & 无穷维概率空间（布朗运动）的存在性无法保证 \\
Borel--Cantelli、强大数定律、0--1 律 & 概率渐近理论全部失效 \\
\bottomrule
\end{tabular}
\end{center}

\textbf{对本笔记的直接影响：} H4 蒙特卡洛积分的收敛性证明依赖 DCT，而 DCT 经 MCT 依赖可数可加。
\textbf{有限可加下 MC 渲染没有理论保证}——这是 CG 侧最实际的理由。

\textit{(2) 不可数可加太强：只剩原子测度。} 若对\emph{任意}指标集可加，则对任意 $X$：
\[
\mu(X) = \sum_{x\in X}\mu(\{x\})
\]
而非负实数族的无序求和收敛\textbf{当且仅当至多可数项非零}。于是 $\mu(\{x\})=0\ (\forall x) \Rightarrow \mu \equiv 0$：
\begin{itemize}
  \item Lebesgue 测度满足 $\lambda(\{x\})=0$，将被迫 $\lambda(\mathbb{R})=0$——荒谬
  \item 正态分布等一切连续型分布被迫退化（C2 整节失效）
  \item A2.5 的分布理论、狄拉克 $\delta$ 建模（点光源）全部无法表述
\end{itemize}
即：不可数可加\textbf{只允许集中在至多可数个原子上的测度}，把连续世界排除在外。

\textit{(3) 可数可加的等价形式：对集合列的连续性。}
\[
A_n \uparrow A \implies \mu(A_n)\to\mu(A), \qquad
A_n \downarrow \varnothing \implies \mu(A_n)\to 0
\]
这是理解其必要性的钥匙——\textbf{分析学的一切建立在极限上，而有限可加对极限一无所知}。

\textit{(4) 代价：Vitali 三难。} 在选择公理下，
\[
\text{可数可加} \;+\; \text{平移不变} \;+\; \text{定义在 }\mathbb{R}\text{ 的全部子集上}
\]
三者\textbf{不可兼得}（Banach 测度靠放弃可数可加才同时做到）。因此我们放弃“全部子集”，
把定义域限制为 $\sigma$-代数，并做完备化。

\[
\boxed{\text{$\sigma$-代数这个看似繁琐的概念，其存在理由正是为可数可加腾出逻辑空间。}}
\]

\textbf{关键测度：}
\begin{itemize}
  \item 勒贝格测度：$\mathbb{R}^n$ 上的“体积”
  \item 计数测度：$\mu(A) = |A|$
  \item 狄拉克测度：$\delta_x(A) = \mathbf{1}_A(x)$（在 CG 中用于点光源建模）
  \item 豪斯多夫测度：用于分形维数（CG 中的分形地形）
\end{itemize}

\subsubsection{$L^p$ 空间}

\[
L^p(X, \mu) = \left\{f : \int |f|^p\,d\mu < \infty\right\} / \sim
\]

其中 $f \sim g$ 当且仅当 $f = g$ 几乎处处。

\begin{itemize}
  \item $L^1$：绝对可积函数（概率密度函数所在空间）
  \item $L^2$：平方可积函数（\textbf{量子力学的状态空间}，信号处理的核心）
  \item $L^\infty$：本质有界函数
\end{itemize}

\textbf{$L^2$ 是希尔伯特空间}，内积为 $\langle f, g \rangle = \int f\bar{g}\,d\mu$。这是连接分析学与量子力学的桥梁。

%--------------------------------------------------------------
\subsection{A2. 泛函分析}

\subsubsection{赋范空间与巴拿赫空间}

\begin{itemize}
  \item \textbf{赋范空间：} 向量空间 + 范数 $\|\cdot\|$（满足正定性、齐次性、三角不等式）
  \item \textbf{巴拿赫空间：} 完备的赋范空间（每个柯西序列都收敛）
  \item 例子：$\mathbb{R}^n$（有限维），$L^p$（无限维），$C[0,1]$（连续函数空间，sup 范数）
\end{itemize}

\subsubsection{希尔伯特空间}

\textbf{定义：} 完备的内积空间。

\textbf{核心性质：}
\begin{itemize}
  \item 内积诱导范数：$\|x\| = \sqrt{\langle x, x \rangle}$
  \item 柯西--施瓦茨不等式：$|\langle x, y \rangle| \leq \|x\|\,\|y\|$
  \item 平行四边形法则：$\|x+y\|^2 + \|x-y\|^2 = 2\|x\|^2 + 2\|y\|^2$
\end{itemize}

\textbf{正交分解定理：} 若 $M$ 是希尔伯特空间 $H$ 的闭子空间，则 $H = M \oplus M^\perp$。

这就是正交投影的严格基础——在 CG 中，最小二乘拟合、法线估计、球谐函数投影都是这个定理的应用。

\subsubsection{正交基与傅里叶分析}

可分希尔伯特空间存在可数正交基 $\{e_n\}$，任何元素可展开：
\[
f = \sum_{n} \langle f, e_n \rangle\, e_n
\]

\textbf{经典正交基：}
\begin{center}
\begin{tabular}{lll}
\toprule
\textbf{空间} & \textbf{正交基} & \textbf{应用} \\
\midrule
$L^2[0, 2\pi]$ & $\{e^{inx}\}$ & 傅里叶级数，信号处理 \\
$L^2(S^2)$ & 球谐函数 $Y_l^m$ & 环境光照、PRT \\
$L^2(\mathbb{R})$ & Hermite 函数 & 量子谐振子 \\
$L^2[0,1]$ & Legendre 多项式 & 数值积分 \\
$L^2(\mathbb{R}^3)$ & 平面波 $e^{i\mathbf{k}\cdot\mathbf{x}}$ & 衍射、散射 \\
\bottomrule
\end{tabular}
\end{center}

\textbf{Parseval 定理：} $\|f\|^2 = \sum |\langle f, e_n \rangle|^2$（能量守恒在频域的表述）

\subsubsection{有界线性算子与谱理论}
\label{subsec:spectral}

\textbf{算子：} $T: H \to H$，线性且有界（$\|Tx\| \leq C\|x\|$）。

\textbf{自伴算子（Hermitian）：} $\langle Tx, y \rangle = \langle x, Ty \rangle$

\textbf{谱定理（有限维）：} 自伴算子可以对角化，特征值全为实数，特征向量构成正交基。

\textbf{谱定理（无限维）：} 自伴算子的谱可能是连续的。例如位置算子 $\hat{x}$ 在 $L^2(\mathbb{R})$ 上的谱是整个实数轴。

\textbf{在 CG 中的应用：}
\begin{itemize}
  \item 拉普拉斯算子 $\Delta$ 的谱 $\to$ 网格上的振动模式 $\to$ 形状分析
  \item 协方差矩阵的特征分解 $\to$ PCA $\to$ 降维、法线估计
\end{itemize}

\textbf{微分算子的无穷维矩阵表示、矩阵指数与 Taylor 公式、若尔当标准型的推广，见~\ref{subsec:infdim-matrix}~节。}

\subsubsection{分布（广义函数）}

\textbf{动机：} 狄拉克 $\delta$ 函数不是普通函数，但在物理中无处不在。

\textbf{定义：} 分布是测试函数空间 $\mathcal{D} = C_c^\infty$ 上的连续线性泛函。
\[
\langle \delta, \varphi \rangle = \varphi(0)
\]

\textbf{关键分布：}
\begin{itemize}
  \item $\delta(x)$：点源（点电荷、点光源）
  \item $\delta'(x)$：偶极子
  \item $\text{p.v.}\frac{1}{x}$：主值分布
\end{itemize}

在 CG 中：点光源的辐射度用 $\delta$ 函数建模，BRDF 中的镜面反射分量包含 $\delta$ 函数。

%--------------------------------------------------------------
\subsection{A3. 微分学的高阶推广：微分形式}

\textbf{0-形式：} 标量函数 $f$ \\
\textbf{1-形式：} $\omega = P\,dx + Q\,dy + R\,dz$（对应向量场） \\
\textbf{2-形式：} $\omega = P\,dy\wedge dz + Q\,dz\wedge dx + R\,dx\wedge dy$（对应通量） \\
\textbf{3-形式：} $\omega = f\,dx\wedge dy\wedge dz$（对应体密度）

\textbf{外微分 $d$：}
\begin{itemize}
  \item $df = \frac{\partial f}{\partial x}dx + \frac{\partial f}{\partial y}dy + \frac{\partial f}{\partial z}dz$（梯度）
  \item $d(\text{1-形式})$ 对应旋度
  \item $d(\text{2-形式})$ 对应散度
  \item $d^2 = 0$（旋度的散度为零，梯度的旋度为零）
\end{itemize}

\textbf{广义 Stokes 定理：}
\[
\boxed{\int_{\partial M} \omega = \int_M d\omega}
\]

这一个公式统一了：
\begin{itemize}
  \item 微积分基本定理（$n=1$）
  \item 格林公式（$n=2$）
  \item 经典 Stokes 定理（$n=2$，三维中）
  \item 散度定理/高斯定理（$n=3$）
\end{itemize}

在 CG 中：离散外微分（DEC）用于网格上的物理模拟、守恒律的离散保持、网格参数化。

%--------------------------------------------------------------
\subsection{A4. Sobolev 空间与弱解}
\label{subsec:sobolev}

\textbf{为什么单列一节：} A1 的 $L^p$、A2 的 Hilbert 空间、P9 的能量法，三条线在这里汇合，
产出的东西是\textbf{整个椭圆 PDE 存在性理论的骨架}。没有它，「NS 方程有解吗」这句话根本没法说清。

\subsubsection{经典导数的不足}

$a$ 处的导数要求 $u$ 在 $a$ 的邻域内有定义且极限存在。可是：

\begin{itemize}
  \item $u(x)=|x|$ 在 $0$ 处无经典导数，但它是很好的函数，其「导数」应当是符号函数；
  \item 激波、间断初值、Leray--Hopf 弱解——这些\textbf{物理上有意义的解根本不光滑}；
  \item 极限过程会破坏光滑性：$C^1$ 中的序列可以在 $L^p$ 中收敛到一个不 $C^1$ 的函数。
\end{itemize}

\textbf{出路：} 用分部积分把导数「搬」到测试函数上——这正是 P9「移导数」那一手，
只不过现在把它\textbf{当作定义}。

\subsubsection{弱导数}

\textbf{定义：} 设 $u\in L^1_{\text{loc}}(\Omega)$。若存在 $v\in L^1_{\text{loc}}(\Omega)$，使对一切
$\varphi\in C_c^\infty(\Omega)$（紧支光滑测试函数）成立
\[
\int_\Omega u\,\partial_i\varphi\,dx=-\int_\Omega v\,\varphi\,dx,
\]
则称 $v$ 是 $u$ 关于 $x_i$ 的\textbf{弱导数}，记 $v=\partial_i u$。

\textbf{直觉：} 经典分部积分 $\int u\varphi'=-\int u'\varphi$（边界项因 $\varphi$ 紧支而消失）本来是
\emph{由 $u'$ 推出}的恒等式；现在反过来，\textbf{用它来定义 $u'$}。$u$ 只需要可积，不需要连续。

弱导数若存在则\textbf{几乎处处唯一}，且与经典导数在后者存在时一致。

\subsubsection{Sobolev 空间}

\[
W^{k,p}(\Omega)=\Big\{u\in L^p(\Omega):\ \partial^\alpha u\in L^p(\Omega)\ \text{对一切}\ |\alpha|\le k\Big\}
\]
（$\alpha$ 是多重指标，$\partial^\alpha$ 按弱导数理解）配备范数
\[
\|u\|_{W^{k,p}}=\Big(\sum_{|\alpha|\le k}\|\partial^\alpha u\|_{L^p}^p\Big)^{1/p}
\qquad(p<\infty)
\]

记 $H^k(\Omega):=W^{k,2}(\Omega)$。$H_0^k(\Omega)$ 是 $C_c^\infty(\Omega)$ 在 $H^k$ 范数下的闭包——
\textbf{「零边界条件」的精确含义}。

\begin{center}
\begin{tabular}{L{2.4cm}L{5.6cm}L{6.2cm}}
\toprule
\textbf{空间} & \textbf{性质} & \textbf{要点} \\
\midrule
$W^{k,p}$ & 完备（Banach） & 所以能取极限、能套压缩映射、能谈存在唯一 \\
$H^k=W^{k,2}$ & Hilbert（内积 $\sum_{|\alpha|\le k}\langle\partial^\alpha u,\partial^\alpha v\rangle$） &
可直接用 A2 的正交分解与投影，$H_0^1$ 更是 \\
磨光（mollification） & $C_c^\infty(\Omega)$ 在 $W^{k,p}(\Omega)$ 中稠密（$p<\infty$） &
\textbf{光滑函数够用}：先在光滑里证，再取极限 \\
\bottomrule
\end{tabular}
\end{center}

\subsubsection{嵌入定理}

\textbf{Sobolev 嵌入：} 设 $\Omega\subset\mathbb R^n$ 有界、边界「足够好」，$1\le p<n$，则
\[
W^{1,p}(\Omega)\hookrightarrow L^{p^*}(\Omega),\qquad
\frac1{p^*}=\frac1p-\frac1n
\]
（$p>n$ 时 $W^{1,p}\hookrightarrow C^{0,\alpha}$，$\alpha=1-n/p$——变成了连续函数。）

\textbf{怎么读这个公式：} 多「一个」导数的可积性，换一个「更好」的空间；
\textbf{损失的额度由维数 $n$ 决定}。$n=1$ 时 $p^*=\infty$：一维情形 $W^{1,p}\hookrightarrow C$，
所以一维「光滑性」几乎是免费的——这也是为什么一维 ODE 理论比 PDE 简单得多。

\textbf{紧嵌入（Rellich--Kondrachov）：} 有界 $\Omega$ 上，$p<n$ 时
\[
W^{1,p}(\Omega)\hookrightarrow\hookrightarrow L^q(\Omega)\qquad(\textstyle 1\le q<p^*)
\]
\textbf{「紧」是这里的全部重点：} 有界集映成相对紧集 $\Longrightarrow$ 可取收敛子列
（见 Q5：这就是变分法直接方法要的那一步）。$q<p^*$ 的严格不等号不能改成等号——
临界指数处嵌入恰好不紧，这正是\textbf{临界指数问题}难的原因。

\textbf{Poincaré 不等式：}
\[
\|u-u_\Omega\|_{L^p}\le C\,\|\nabla u\|_{L^p}\qquad
\text{特别地}\quad \|u\|_{L^p}\le C\|\nabla u\|_{L^p}\ \ (u\in W_0^{1,p})
\]
后一式说明在 $H_0^1$ 上 $\|\nabla u\|_{L^2}$ \textbf{本身就是等价范数}——
它是下面「强制性」的来源。

\subsubsection{弱解与 Lax--Milgram}

以 Poisson 方程为例：$-\Delta u=f$，$u|_{\partial\Omega}=0$。

\textbf{弱形式：} 两边乘 $v\in H_0^1$、积分、对左边分部积分（P9 的移导数），求 $u\in H_0^1$ 使
\[
\underbrace{\int_\Omega\nabla u\cdot\nabla v\,dx}_{B(u,v)}=\underbrace{\int_\Omega fv\,dx}_{F(v)}
\qquad\forall v\in H_0^1
\]

\textbf{Lax--Milgram 定理：} $H$ 是 Hilbert 空间，$B:H\times H\to\mathbb R$ 双线性。若
\begin{itemize}
  \item \textbf{有界：} $|B(u,v)|\le C\|u\|\|v\|$，
  \item \textbf{强制（coercive）：} $B(u,u)\ge\alpha\|u\|^2$（$\alpha>0$），
\end{itemize}
且 $F$ 是连续线性泛函，则存在\textbf{唯一} $u\in H$ 使 $B(u,v)=F(v)$ 对一切 $v$ 成立。

\textbf{验证只需两步（这就是它的价值）：}
\begin{itemize}
  \item 有界：$|B(u,v)|\le\|\nabla u\|\|\nabla v\|$（直接 Cauchy--Schwarz，P3 的第一件套）；
  \item 强制：$B(u,u)=\|\nabla u\|^2\ge\alpha\|u\|_{H^1}^2$——\textbf{正是 Poincaré 不等式}。
\end{itemize}

\textbf{方法论上的意义：} 它把「解一个无穷维方程」变成「验证两个不等式」。
弱解存在唯一性一旦拿到，剩下的就是\textbf{正则性}问题（把弱解升级成经典解）。

\subsubsection{正则性提升与能量估计}

\textbf{能量估计（先验估计）：} 在弱形式中取 $v=u$，立刻得
\[
\|\nabla u\|_{L^2}^2=\int f u\le\|f\|_{L^2}\|u\|_{L^2}\le C\|f\|_{L^2}\|\nabla u\|_{L^2}
\ \Longrightarrow\ \|u\|_{H^1}\le C\|f\|_{L^2}
\]
（中间一步用 Cauchy--Schwarz，最后一步用 Poincaré。）\textbf{这不是「顺便得到」的结论——
PDE 理论中几乎所有存在性都是先拿到先验估计，再用紧性取极限。}

\textbf{内部正则性：} 若 $f\in H^k$ 且 $\partial\Omega$ 光滑，则 $u\in H^{k+2}$，
$k$ 足够大时由 Sobolev 嵌入即得 $u\in C^{2,\alpha}$——\textbf{弱解自动变成经典解}。

\textbf{这条链要记住：}
\[
\text{弱形式}\to\text{Lax--Milgram 存在唯一}\to\text{能量估计}\to\text{正则性提升}\to\text{经典解}
\]

\textbf{连到 NS：} 同一条链在 Navier--Stokes 上只能走到第二步——
Leray--Hopf 弱解的存在性（1934）靠的正是 P9 的能量法流水线（乘 $u$、移导数、Gronwall 收尾）。
\textbf{但正则性与唯一性至今未解决}（千禧年问题）。所以 PDE 里「弱解」不是权宜之计，
而是\textbf{非线性和临界情形下唯一能拿到的东西}。

\textbf{与既有章节的接口：}
\begin{itemize}
  \item $L^p$ 与收敛定理见 A1；Hilbert 空间与投影见 A2；
  \item 移导数与能量法五步见 P9；紧嵌入为什么关键见 Q5；Rellich--Kondrachov 的拓扑背景见 Q5 末；
  \item 分布（广义函数）是弱导数的更一般版本，见 A2.5。
\end{itemize}